\chapter{Introduction}
\label{ch:introduction}

\section{Motivation and Background}
\label{sec:motivation}

The past decade has witnessed an unprecedented surge in data-parallel workloads driven
by machine learning inference, computer vision, digital signal processing, and scientific
simulation. Modern neural-network inference passes millions of multiply-accumulate
operations over tensors that may contain thousands of elements of identical type.
A classical scalar processor that fetches, decodes, and executes one operation per
clock cycle cannot amortise the instruction-fetch overhead across these repeated
computations; it spends the majority of its energy on control flow rather than useful
arithmetic. The fundamental insight that motivates vector processing is simple: if the
same arithmetic operation is to be applied to a long sequence of operands, a single
instruction can describe the entire sequence, and the hardware can overlap the execution
of individual elements in time, in space, or both.

Vector processing architectures have a history stretching back to the Cray-1 vector
supercomputer of 1976, yet they have undergone a dramatic renaissance in the era of
deep learning. Modern CPU vector extensions --- Intel AVX-512, ARM SVE/SVE2, and
RISC-V RVV --- share the same conceptual heritage: a set of wide registers holding
multiple data elements and an instruction vocabulary that operates on all elements
simultaneously \cite{simd-intro}. These extensions achieve substantial performance
improvements on workloads such as image convolution, matrix multiplication, and
activation functions, often by factors of four to sixteen compared to scalar code on
the same core, with minimal additional area overhead when the vector width is chosen
judiciously.

The RISC-V open instruction set architecture has emerged as a compelling foundation for
custom silicon development \cite{riscv-spec-2019, patterson2020risc}. Its permissive
licensing model, its separation of mandatory base ISA from optional standard extensions,
and its active ecosystem of open-source toolchains and verification infrastructure make
it uniquely attractive for academic and start-up design teams that cannot afford
proprietary ISA licenses. The RISC-V Vector Extension (RVV), standardised as version
1.0 in 2021 \cite{rvv-spec-1-0}, defines a length-agnostic programming model that
allows the same binary to run correctly on vector units ranging from 128-bit to 65536-bit
register widths, simply by rerunning the \texttt{vsetvli} configuration instruction
at runtime. This future-proof abstraction is particularly valuable for \gls{asic}
design because the hardware can be dimensioned to the minimum area required for the
target workload without invalidating software compiled for wider implementations.

The choice between \gls{fpga} prototyping and \gls{asic} tapeout reflects a fundamental
engineering tradeoff. An \gls{fpga} provides rapid iteration and in-system debugging
at the cost of power efficiency, achievable clock frequency, and configurable logic
resources. An \gls{asic}, once fabricated, consumes three to five times less power and
operates at two to four times higher frequency for equivalent logic, because every
transistor is sized and placed for a single purpose \cite{weste2010cmos}. For embedded
inference at the edge --- smart cameras, wearable health monitors, autonomous vehicle
perception nodes --- the power envelope is critical, and \gls{asic} implementation is
the only viable path. Producing tapeout-ready \gls{rtl} therefore has direct industrial
relevance beyond the academic context of this thesis.

Writing \gls{rtl} that is not merely functionally correct in simulation but also fully
synthesisable, timing-clean, and free of latches and X-propagation paths is a
substantially harder problem than simulation-only design. Synthesisers reject
\texttt{initial} blocks, non-synthesisable system tasks, and asynchronous latches.
Static timing analysis tools expose combinational paths that violate setup and hold
constraints. Any path that is not properly constrained can silently cause functional
failures after place-and-route. This thesis takes the position that a vector processor
should be designed to these tighter standards from the outset, rather than retrofitted
after the fact.

\section{Problem Statement}
\label{sec:problem}

Designing a standards-compliant vector processor presents multiple interacting
technical challenges that make it substantially more complex than a straightforward
scalar core.

\paragraph{ISA Complexity.}
The RVV Zve32x subset encompasses more than fifty distinct instruction encodings across
arithmetic, logical, shift, compare, reduction, and load/store categories. Each
instruction must be correctly decoded from a 32-bit encoding that shares the OP-V
opcode but varies along three independent axes: \texttt{funct6} (6 bits specifying
the operation), \texttt{funct3} (3 bits selecting the operand type: VV, VX, or VI),
and \texttt{vm} (1 bit enabling or disabling masking). The decoder must distil this
information into a compact control word that drives downstream execution units over
potentially many clock cycles, because vector operations execute once per element rather
than once per instruction.

\paragraph{Hazard Management.}
A combined scalar--vector system requires hazard resolution at two distinct time
scales. The scalar pipeline must detect read-after-write hazards between consecutive
instructions on a cycle-by-cycle basis using forwarding and load-use stalling.
Concurrently, the scalar core must stall when it encounters a vector instruction while
the \gls{vpu} instruction queue is full, and it must stall again when a
\texttt{vsetvl}/\texttt{vsetvli} instruction requires the \gls{vpu} to complete its
configuration and write the computed vector length back to a scalar register. Both
stall conditions must be detected and propagated through the pipeline without corrupting
any in-flight instructions.

\paragraph{Area and Power Tradeoffs.}
A multi-lane \gls{vpu} with VLEN=128 bits and four 32-bit lanes executing simultaneously
yields peak throughput but occupies proportionally more area. For an ASIC tapeout
constrained by educational budget, the single-lane implementation trades peak throughput
for area efficiency, completing one 32-bit element per cycle. The reduction unit is
particularly sensitive: a na\"ive reduction tree with sixteen adders connected in a
combinational chain produces a long critical path that limits \texttt{Fmax}. The
design must replace this with a serialised approach that achieves correct results while
remaining timing-clean.

\paragraph{Masked Operations.}
RVV masking allows each vector element to be conditionally suppressed based on the
contents of vector register \texttt{v0}. For compare instructions that \emph{produce}
a mask result, the implementation must accumulate per-element bits into a mask buffer
and commit the result atomically at the end of the instruction, rather than writing
partial results to the register file mid-execution. This requires an additional FSM
state (\texttt{ST\_FINAL\_MASKING}) and a separate mask write buffer, both of which
add control complexity without adding arithmetic throughput.

\paragraph{Memory Interface.}
A Harvard memory architecture with separate instruction and data memories simplifies
timing by eliminating structural hazards on the memory bus, but requires careful
arbitration when both the scalar \gls{lsu} and the vector \gls{vlsu} wish to access
data memory simultaneously. The dual-port memory solution must be synthesisable (i.e.,
mappable to block-RAM primitives) and must not introduce dead cycles into either
the scalar or vector memory access path.

\section{Objectives}
\label{sec:objectives}

This project has the following specific technical objectives:

\begin{enumerate}[label=\arabic*.]
  \item \textbf{Implement the \texttt{rv32im\_zicsr\_zve32x\_zvl128b} ISA.}
        Deliver a scalar core supporting the RV32I base, the M multiplication/division
        extension, and the Zicsr privileged instruction extension, integrated with a
        vector coprocessor that implements the full Zve32x minimum vector subset at
        VLEN=128 bits.

  \item \textbf{Five-stage pipelined scalar core.}
        The scalar core shall implement IF, ID, EX, MEM, and WB stages with
        EX$\rightarrow$MEM$\rightarrow$WB data forwarding and load-use hazard stalling,
        achieving CPI approaching 1.0 on non-memory, non-vector code.

  \item \textbf{Single-lane, multi-cycle VPU with full Zve32x coverage.}
        The \gls{vpu} shall process one SEW-wide element per clock cycle via a single
        execution lane containing an adder, multiplier, shifter, logic, comparator,
        min/max, and reduction sub-unit. It shall support SEW=8/16/32 and
        LMUL=1/2/4/8.

  \item \textbf{Tapeout-viable RTL.}
        All modules in the \texttt{/rtl/} hierarchy shall be synthesisable: no
        \texttt{initial} blocks, no \texttt{\$display} in synthesisable paths, no
        combinational latches, synchronous active-low reset throughout, and all widths
        parameterised to avoid hardcoded VLEN/SEW bit slices.

  \item \textbf{Comprehensive functional verification.}
        A directed testbench suite shall achieve 172/172 test cases passing before any
        design change is considered complete. Integration benchmarks shall include
        BT.601 RGB-to-grayscale conversion on a 128$\times$128 image and general
        matrix multiplication, both verified at pixel/element level.

  \item \textbf{UART-based end-to-end system test.}
        A UART peripheral mapped at \texttt{0xFF000000} shall enable real-data
        streaming from a host computer to the running system. Firmware shall receive
        raw image data, invoke the vector grayscale pipeline, and return the luminance
        channel, providing an end-to-end integration test beyond pure simulation.
\end{enumerate}

\section{Contributions}
\label{sec:contributions}

The following are the original contributions of this thesis:

\begin{enumerate}[label=\arabic*.]
  \item \textbf{A complete open-source Zve32x implementation in SystemVerilog.}
        This work delivers the first publicly documented single-lane Zve32x
        implementation that explicitly targets ASIC tapeout rules, including
        parameterised widths, synchronous active-low reset, no latches, and no
        simulation-only constructs in any production RTL file.

  \item \textbf{A 48-bit structured control bus architecture.}
        The \texttt{vproc\_vdecoder} module produces a 48-bit control word that
        encodes all execution context --- register addresses, operation flags,
        masking controls, reduction flags, and vtype --- into a single bus that passes
        through the instruction \gls{fifo} and drives all downstream units without
        any re-decoding. This single-decode, single-pipeline-register design minimises
        routing complexity and ensures synthesis tools can easily register the control
        path.

  \item \textbf{A serial reduction unit with busy-based protocol.}
        Unlike published designs that employ a combinational reduction tree (which
        produces a $O(\log_2 N)$-deep critical path), the \texttt{vproc\_reduction}
        module accumulates one element per clock cycle using a single adder under
        FSM control. The \texttt{reduction\_busy} signal allows the FSM to overlap
        the next instruction's decode with the final accumulation cycle, improving
        utilisation without increasing logic depth.

  \item \textbf{A dual-port DMEM arbitration scheme for scalar--vector coexistence.}
        The \texttt{dmem\_sync} module exposes independent scalar and VLSU ports that
        map to the two read/write ports of a single block-RAM primitive, eliminating
        structural hazards between scalar load/store and vector load/store without
        an explicit arbiter.

  \item \textbf{A UART TL-UL slave interface for end-to-end image streaming.}
        A synthesisable UART peripheral following the TileLink Uncached Lightweight
        (\gls{tlu}) protocol is integrated at the top-level address space.
        Firmware running on the RISC-V core can stream raw image data bytes via
        MMIO writes and reads, enabling end-to-end validation that no simulation-level
        abstraction could provide.

  \item \textbf{Validated FPGA prototype on DE10-Standard achieving 47.18~MHz.}
        The complete design was synthesised and placed-and-routed using Intel Quartus
        Prime on a Cyclone~V device, consuming 10\,438 \glspl{alm} and achieving an
        \gls{fmax} of 47.18~MHz --- a $3\times$ frequency improvement over an earlier
        iteration that used a combinational reduction chain, demonstrating the benefit
        of the serial reduction redesign.
\end{enumerate}

\section{Thesis Organisation}
\label{sec:organisation}

The remainder of this thesis is organised as follows.

\textbf{Chapter~\ref{ch:background}} provides the theoretical and contextual background
required to understand the design decisions made in subsequent chapters. It surveys
the RISC-V instruction set architecture, the principles of vector processing and how
they differ from scalar SIMD, the RVV specification in detail, and the standard ASIC
design methodology from \gls{rtl} through tapeout. The chapter concludes with a
comparative survey of related academic and industrial vector processor implementations,
positioning this work within the broader landscape.

\textbf{Chapter~\ref{ch:specification}} defines the complete system specification that
guided the implementation. It states the functional and non-functional requirements,
provides a detailed instruction table covering all supported encodings, specifies the
vector configuration model (VLEN, SEW, LMUL, VLMAX), describes the Harvard memory
architecture with precise address maps, specifies the UART peripheral's register file
and timing constraints, and states the performance targets against which the design
is evaluated.

\textbf{Chapter~\ref{ch:hardware}} describes the hardware architecture and
\gls{rtl} implementation in detail. It covers the five-stage scalar pipeline with
hazard detection and forwarding, the vector decoder and its 48-bit control bus, the
instruction \gls{fifo}, the nine-state execution \gls{fsm}, the single-lane execution
pipeline and each of its sub-units, the vector register file, the vector load/store unit,
and the UART peripheral.

\textbf{Chapter~\ref{ch:software}} describes the firmware and software infrastructure.
It covers the assembly language programming model exposed to firmware writers, the
build system and toolchain configuration, the BT.601 grayscale benchmark, the UART
streaming protocol, and the \texttt{elf2hex} conversion pipeline that loads programs
into the hardware.

\textbf{Chapter~\ref{ch:evaluation}} presents verification and evaluation results.
It describes the testbench methodology, the 172-case regression suite, per-instruction
pass/fail analysis, the 128$\times$128 pixel grayscale benchmark, FPGA timing and
area results, and a discussion of the frequency bottleneck and how the serial reduction
redesign resolved it.

\textbf{Chapter~\ref{ch:conclusion}} summarises the contributions, discusses limitations
of the current implementation, and outlines directions for future work including
multi-lane scaling, floating-point support, and formal verification.
